\subsection{Regresión por mínimos cuadrados ordinarios: un paso}
\label{subsec:ols-univariante}

La versión actual de \texttt{root/OLS.py} recibe únicamente generación eléctrica y predice generación eléctrica. La fecha y hora se utilizan para seleccionar el periodo, ordenar las observaciones, comprobar el intervalo entre registros e identificar las predicciones; no son variables explicativas.

Se lee \texttt{datos.csv} y se interpreta \texttt{Timestamp} conservando horas y minutos. Se selecciona la primera área distinta de \texttt{Area}, desde el 1 de enero de 2022, inclusive, hasta el 1 de enero de 2023, exclusivo. Este intervalo incluye el 31 de diciembre completo. Después del ordenamiento cronológico, se conserva el primer 70\,\% de los registros; el 30\,\% final no se evalúa.

El código comprueba que la generación sea numérica y finita, que no existan fechas duplicadas y que las diferencias entre fechas consecutivas sean constantes. Si estas condiciones no se cumplen, detiene la ejecución. Una historia de 24 registros equivale a 24 horas únicamente cuando el intervalo observado es de una hora.

Si $g_t$ representa la generación del registro $t$, cada ejemplo se define mediante
\begin{equation}
\boldsymbol{x}_t=(g_{t-24},g_{t-23},\ldots,g_{t-1})^{\mathsf T},
\qquad y_t=g_t.
\end{equation}
Se construye una matriz con 24 columnas de rezagos. Para cada predicción se utilizan los 24 ejemplos supervisados inmediatamente anteriores: \texttt{X\_train} tiene dimensión $24\times24$, \texttt{y\_train} contiene 24 generaciones y \texttt{X\_test} tiene dimensión $1\times24$. Los ejemplos son solapados y comparten parte de su historia.

Para predecir $g_t$ se ajusta el modelo con los pares $(\boldsymbol{x}_s,g_s)$ para $s=t-24,\ldots,t-1$. Estos ejemplos requieren observaciones desde $g_{t-48}$, por lo que la primera predicción evaluada es la del registro 49. El programa exige al menos 50 registros dentro del primer 70\,\% para evaluar dos o más predicciones.

En cada ajuste, \texttt{StandardScaler} se estima exclusivamente con las entradas de entrenamiento y se aplica a la entrada de prueba. La respuesta conserva su escala en MWh. Se utiliza \texttt{LinearRegression} para estimar el intercepto y los 24 coeficientes:
\begin{equation}
(\hat{\beta}_{0,t},\hat{\boldsymbol{\beta}}_t)
=\arg\min_{\beta_0,\boldsymbol{\beta}}
\sum_{s=t-24}^{t-1}
\left(g_s-\beta_0-\widetilde{\boldsymbol{x}}_s^{\mathsf T}
\boldsymbol{\beta}\right)^2,
\end{equation}
donde $\widetilde{\boldsymbol{x}}_s$ contiene los rezagos estandarizados con los parámetros del ajuste actual. La salida es
\begin{equation}
\hat g_t=\hat{\beta}_{0,t}
+\widetilde{\boldsymbol{x}}_t^{\mathsf T}\hat{\boldsymbol{\beta}}_t.
\end{equation}

La ventana avanza un registro y se repiten el escalamiento y el ajuste. Para el siguiente origen se usan las generaciones observadas disponibles hasta ese instante. No se encadenan las predicciones anteriores como entradas: es una evaluación de un paso con reajuste móvil.

\begin{table}[htbp]
\centering
\caption{Desplazamiento de las ventanas del OLS univariante. Los índices identifican registros ordenados.}
\label{tab:ventanas-ols}
\begin{tabular}{lll}
\toprule
Evaluación & Respuestas de entrenamiento & Entrada de prueba y respuesta \\
\midrule
1 & $g_{25},\ldots,g_{48}$ & $(g_{25},\ldots,g_{48})\to g_{49}$ \\
2 & $g_{26},\ldots,g_{49}$ & $(g_{26},\ldots,g_{49})\to g_{50}$ \\
3 & $g_{27},\ldots,g_{50}$ & $(g_{27},\ldots,g_{50})\to g_{51}$ \\
\bottomrule
\end{tabular}
\end{table}

Se calculan MAE, RMSE y $R^2$, y se muestran el área, el intervalo temporal, las fechas evaluadas y el número de predicciones. La gráfica usa fecha y hora en el eje horizontal. Con 24 ejemplos y 24 predictores más intercepto, los coeficientes pueden no ser únicos; la interpretación debe apoyarse en la evaluación de observaciones posteriores al ajuste.
